\documentclass[11pt, letterpaper]{article}
\usepackage{amsmath, amssymb}
\usepackage{hyperref}
\usepackage{booktabs}
\usepackage{enumitem}
\usepackage{xcolor}
\usepackage{geometry}
\usepackage{fancyhdr}
\usepackage{titlesec}
\usepackage{tcolorbox}
\usepackage{parskip}
\geometry{margin=1in}

\definecolor{courseblue}{HTML}{002452}
\definecolor{coursegold}{HTML}{FABD0F}
\definecolor{coursered}{HTML}{B90E31}
\definecolor{lightblue}{HTML}{E8EEF7}

\hypersetup{colorlinks=true,linkcolor=courseblue,urlcolor=coursered}

\tcbuselibrary{skins}
\newtcolorbox{guidancebox}{enhanced,colback=lightblue,colframe=courseblue,
  sharp corners,leftrule=4pt,left=6pt,right=6pt,top=4pt,bottom=4pt}

\pagestyle{fancy}\fancyhf{}
\fancyhead[L]{\small\textcolor{courseblue}{\textbf{CISC 473 --- Deep Learning | Fall 2026}}}
\fancyhead[R]{\small\textcolor{courseblue}{Project Proposal}}
\fancyfoot[C]{\small\thepage}
\renewcommand{\headrulewidth}{0.4pt}

\titleformat{\section}{\large\bfseries\color{courseblue}}{}{0em}{}[\titlerule]
\titleformat{\subsection}{\normalsize\bfseries\color{courseblue}}{}{0em}{}

%% ════════════════════════════════════════════════════════════════
%% INSTRUCTIONS FOR STUDENTS
%% Replace every [PLACEHOLDER] with your actual content.
%% Delete all \guidancebox blocks before submitting.
%% Submit as a compiled PDF + .tex source in a .zip archive.
%% ════════════════════════════════════════════════════════════════

\begin{document}

{\centering
  {\Large\textbf{CISC 473: Deep Learning --- Project Proposal}}\\[0.5em]
  {\large\textit{P08: Abstractive Summarization via Parameter-Efficient Fine-Tuning}}\\[0.4em]
  \textbf{Group Members:}\\{}
  Choo Wen Bin \quad 20630085 \quad (Role: Project Lead)\\{}
  Nai Junn Rong \quad 20630086 \quad (Role: Research Lead)\\{}

  \textbf{GitHub Repository:} \url{https://github.com/wenbin05/DeepLearningP08}\\[0.3em]
  \textbf{Date:} \today\\[0.4em]
  \rule{\linewidth}{0.4pt}
}

\section{Problem Statement and Motivation}


Abstractive text summarization aims to generate concise and coherent summaries
that capture the important information in a source document using newly
generated language, rather than simply extracting sentences from the original
text. This task is challenging because a good summary must simultaneously
identify salient information, preserve the meaning of the source, remain
linguistically fluent, and avoid introducing unsupported or contradictory
facts. Modern encoder--decoder Transformer models such as BART have achieved
strong performance on abstractive summarization by leveraging large-scale
pre-training before being adapted to downstream datasets~[1].

However, adapting large pre-trained models through conventional full
fine-tuning requires updating most or all of their parameters. As model sizes
increase, this creates substantial computational, memory, and storage costs,
especially when separate versions of the same model must be maintained for
different tasks or domains. Parameter-efficient fine-tuning (PEFT) methods
address this problem by keeping most of the pre-trained model fixed and
training only a small number of additional or modified parameters. Methods such
as Low-Rank Adaptation (LoRA)~[2] and Prefix Tuning~[3] have therefore become
important approaches for reducing the cost of adapting pre-trained Transformer models.

Despite their efficiency, an important question remains: how should limited
adaptation capacity be used while preserving high-quality summarization performance?
Reducing the number of trainable parameters may improve efficiency, but it can
also affect summary quality, factual consistency, and the ability to generalize
to text from a different domain. These trade-offs are important in practical
applications such as news summarization, document analysis, and information
retrieval, where models may need to be adapted repeatedly under limited
computational resources. This project therefore studies parameter-efficient
adaptation of BART for abstractive summarization, with particular attention to
the balance between efficiency, summarization quality, factual reliability,
and cross-domain generalization.

\section{Literature Survey}


\subsection{Related Works}

\paragraph{BART: Lewis et al., 2020~[1].}
Lewis et al.\ introduced BART, a denoising sequence-to-sequence pre-training
framework that combines a bidirectional Transformer encoder with an
autoregressive decoder. The model is pre-trained by corrupting text and
learning to reconstruct the original sequence, allowing it to perform strongly
on conditional generation tasks such as abstractive summarization. On the
CNN/DailyMail benchmark, BART-large achieved ROUGE-1/2/L scores of
44.16/21.28/40.90, while on XSum it achieved 45.14/22.27/37.25~[1].
These results established BART as a strong summarization architecture and
motivate its use as the backbone of our project. However, Lewis et al.\
adapted BART through conventional full fine-tuning, whereas our project
investigates whether comparable summarization performance can be obtained by
training only a small fraction of the model parameters.

\paragraph{LoRA: Hu et al., 2022~[2].}
Hu et al.\ proposed Low-Rank Adaptation (LoRA), which freezes a pre-trained
weight matrix $W_0$ and represents its task-specific update using two much
smaller trainable matrices such that $\Delta W = BA$. This reduces the number
of parameters that must be optimized while allowing the learned update to be
merged into the original model at inference time. On GPT-3 175B, a LoRA
configuration with 4.7 million trainable parameters achieved
ROUGE-1/2/L scores of 53.8/29.8/45.9 on the SAMSum summarization dataset,
compared with 52.0/28.0/44.5 for full fine-tuning of approximately
175 billion parameters~[2]. LoRA is the main PEFT method studied in our
project. In contrast to the original paper's broad evaluation across several
models and tasks, we focus specifically on BART-based abstractive
summarization and study how LoRA rank and placement influence performance,
efficiency, and transfer.

\paragraph{Prefix-Tuning: Li and Liang, 2021~[3].}
Li and Liang proposed Prefix-Tuning, a parameter-efficient alternative to
full fine-tuning in which the pre-trained model remains frozen while a small
set of continuous task-specific prefix vectors is optimized. These learned
vectors act as virtual tokens that subsequent Transformer activations can
attend to. Using BART-large on XSum, full fine-tuning achieved
ROUGE-1/2/L scores of 45.14/22.27/37.25, while Prefix-Tuning achieved
43.80/20.93/36.05 using a 2\% configuration and
42.92/20.03/35.05 using approximately 0.1\% task-specific parameters~[3].
The authors also found stronger extrapolation to unseen topics in some
settings. Prefix-Tuning therefore provides an important PEFT baseline in our
project because it adapts the model through learned conditioning vectors,
rather than through low-rank modifications of its weight matrices as in LoRA.

\paragraph{AdaLoRA: Zhang et al., 2023~[4].}
Zhang et al.\ observed that assigning the same low-rank adaptation budget to
every weight matrix implicitly assumes that all model components are equally
important, and proposed AdaLoRA to allocate adaptation capacity dynamically
according to parameter importance. AdaLoRA uses an SVD-inspired
parameterization and progressively removes less important rank components
while preserving more capacity in influential matrices. With BART-large and a
1.10\% trainable-parameter budget, AdaLoRA achieved ROUGE-1/2/L scores of
44.35/21.13/36.13 on XSum, compared with 43.40/20.20/35.20 for standard
LoRA. On CNN/DailyMail, the corresponding scores were
44.96/21.77/42.09 and 44.72/21.58/41.84~[4]. This work is closely related
to our project because it demonstrates that the allocation of a fixed
adaptation budget can affect performance. Our study differs by examining
explicit and interpretable architectural locations within BART rather than
learning the allocation dynamically.

\paragraph{WeightLoRA: Veprikov et al., 2026~[5].}
Veprikov et al.\ proposed WeightLoRA, which treats individual LoRA adapters
as selectable units and learns which adapters are sufficiently important to
retain during fine-tuning. Their WeightLoRA+ variant reallocates the saved
parameter budget by increasing the rank of selected adapters. In BART-large
experiments with rank $r=8$, standard LoRA used approximately 1.05\% of the
model parameters and obtained ROUGE-1 scores of 36.07 on XSum and 25.28 on
CNN/DailyMail, whereas WeightLoRA+ used the same reported parameter percentage
and obtained 36.74 and 25.72, respectively~[5]. These results provide recent
evidence that the usefulness of LoRA adapters is not uniform throughout a
model. Our project differs from WeightLoRA because we do not learn which
individual adapters to retain; instead, we perform a controlled comparison of
architecturally distinct BART regions under approximately matched trainable
parameter budgets and additionally examine cross-domain transfer and factual
consistency.


\paragraph{FactCC: Kryscinski et al., 2020~[6].}
Kryscinski et al.\ argued that lexical-overlap metrics such as ROUGE do not
directly determine whether a generated summary is factually consistent with
its source and introduced FactCC, a weakly supervised factual-consistency
classifier for abstractive summarization. The method generates synthetic
consistent and inconsistent claims from source documents and trains a
BERT-based classifier to determine whether a claim is supported by the
source. FactCC achieved 72.65\% class-balanced accuracy and a weighted F1 score of
0.86 on its annotated CNN/DailyMail evaluation set, but its balanced
accuracy decreased to 54.11\% on XSum~[6]. This result is relevant to our
project because it shows both the importance and difficulty of factuality
evaluation across domains. We therefore evaluate factual consistency in
addition to ROUGE and retain human evaluation rather than relying exclusively
on an automatic factuality metric.


Taken together, these studies establish the progression motivating our work.
BART provides a strong encoder--decoder architecture for abstractive
summarization, while LoRA and Prefix Tuning demonstrate that downstream
adaptation can be performed with far fewer trainable parameters than full
fine-tuning. AdaLoRA and WeightLoRA further indicate that adaptation capacity
is not equally useful across all model parameters or adapters. At the same
time, FactCC shows that factual consistency cannot be assessed solely through
lexical-overlap metrics such as ROUGE. These findings motivate a controlled
study of where a limited LoRA budget should be placed within BART and whether
that choice affects not only ROUGE performance but also cross-domain
generalization and factual consistency.

\subsection{Open Problem}

AdaLoRA~[4] and WeightLoRA~[5] indicate that adaptation capacity need not be
distributed uniformly across Transformer parameters. However, these works do
not provide a controlled comparison, in BART-based abstractive summarization,
of encoder self-attention, decoder self-attention, and encoder--decoder
cross-attention under approximately matched trainable-parameter budgets.

This distinction is important because these components perform different roles
in an encoder--decoder model. A direct comparison without controlling parameter
count would be confounded: adapting more modules generally introduces more
trainable parameters, making it unclear whether a performance improvement is
caused by the architectural location of adaptation or simply by additional
capacity.

Our project addresses this gap by comparing LoRA adaptation across these
architectural regions while keeping their trainable-parameter counts
approximately matched. We will investigate whether the placement of a limited
LoRA budget affects in-domain summarization quality, zero-shot transfer from
CNN/DailyMail to XSum, factual consistency, and training efficiency.

\section{Proposed Methodology}

\subsection{Baseline}

We will use the pre-trained \texttt{facebook/bart-base} checkpoint from
HuggingFace Transformers as our backbone model and the standard
CNN/DailyMail train, validation, and test splits as the primary summarization
benchmark. A zero-shot BART-base model will first be evaluated to establish
pre-adaptation performance. We will then implement the three required
adaptation baselines: full fine-tuning, Prefix Tuning, and Low-Rank Adaptation
(LoRA), using HuggingFace Transformers and the PEFT library.

For LoRA, we will evaluate ranks $r \in \{4,8,16\}$ and perform the required
target-projection ablation over $Q$, $K$, $V$, and all selected attention
projections. Summarization quality will be measured using ROUGE-1, ROUGE-2,
and ROUGE-L. Training configurations will use the same preprocessing,
generation settings, and data splits wherever possible so that differences
between adaptation methods can be compared fairly.

\subsection{Proposed Extension}

Our extension investigates whether the \emph{architectural location} of LoRA
adaptation affects summarization performance when trainable capacity is held
approximately constant. We will construct three placement variants in which
LoRA is applied only to (1) encoder self-attention, (2) decoder self-attention,
or (3) encoder--decoder cross-attention. The same target projections, LoRA
rank, and training hyperparameters will be used across these variants.
Because BART-base contains the same number of layers and hidden dimensionality
in these corresponding attention blocks, this design yields approximately
matched trainable-parameter counts and isolates placement as the main
experimental variable.

We chose this extension because encoder self-attention, decoder self-attention,
and cross-attention perform different roles: representing the source,
modeling generated text, and conditioning generation on the source,
respectively. Our hypothesis is that, even under similar parameter budgets,
LoRA placement will lead to measurable differences in summarization quality
and cross-domain robustness. This would indicate that parameter efficiency
depends not only on how much adaptation capacity is used, but also on where
that capacity is introduced.

\subsection{Evaluation Plan}

Models will be trained on CNN/DailyMail and evaluated on its held-out test set.
To assess cross-domain generalization, the same trained models will also be
evaluated on the standard XSum test split without further fine-tuning.

We will report ROUGE-1, ROUGE-2, and ROUGE-L for summarization quality,
together with the number and percentage of trainable parameters and training
time for efficiency. Factual consistency will additionally be evaluated using
the required human study of 10 participants over 20 generated summaries.
Comparisons will include zero-shot BART-base, full fine-tuning, Prefix Tuning,
standard LoRA, and our three LoRA placement variants. Main configurations will
be repeated with at least three random seeds, and results will be reported as
mean $\pm$ standard deviation.

\subsection{Early Novel Contribution (Optional Bonus)}

Our proposed novel contribution is a budget-matched ablation of LoRA placement
across BART's encoder self-attention, decoder self-attention, and
encoder--decoder cross-attention. By fixing the target projections, LoRA rank,
and training settings, we isolate architectural placement as the main variable
rather than confounding it with trainable capacity. We will test whether
placement affects in-domain ROUGE performance, zero-shot CNN/DailyMail-to-XSum
transfer, factual consistency, and training efficiency.

\section{Project Timeline and Roles}

\begin{center}
\small
\begin{tabular}{lp{7cm}l}
\toprule
\textbf{Week} & \textbf{Tasks} & \textbf{Responsible} \\
\midrule

1 &
Read BART, LoRA, Prefix-Tuning, AdaLoRA, WeightLoRA, and FactCC;
set up GitHub repository and Python environment; load CNN/DailyMail and XSum. &
Both \\

2 &
Finalize research question, placement experiment, and evaluation protocol;
implement dataset preprocessing and ROUGE evaluation pipeline;
complete proposal. \textbf{(Proposal due)} &
Both \\

3 &
Evaluate zero-shot BART-base on CNN/DailyMail; verify generation and ROUGE
pipeline; implement parameter-counting and experiment logging utilities. &
Wen Bin / Junn Rong \\

4 &
Implement LoRA using HuggingFace PEFT; run rank experiments
$r\in\{4,8,16\}$ and begin $Q/K/V$/all target-projection ablations. &
Wen Bin / Junn Rong \\

5 &
Complete LoRA target-projection ablations; implement and evaluate Prefix
Tuning and full fine-tuning baselines under comparable training budgets. &
Junn Rong / Wen Bin \\

6 &
Evaluate trained models on XSum without further fine-tuning; compile
preliminary quality, efficiency, and cross-domain results; perform initial
analysis and write midterm report. \textbf{(Midterm due)} &
Both \\

7 &
Implement encoder-only, decoder-only, and cross-attention LoRA placement
variants; verify approximately matched trainable-parameter counts; prepare and
conduct the human factual-consistency evaluation. &
Wen Bin / Junn Rong \\

8 &
Complete placement experiments with multiple random seeds; analyse human
factuality results; measure training time and parameter efficiency; conduct
extended ablations where required. &
Both \\

9 &
Aggregate mean $\pm$ standard deviation results; produce ROUGE-versus-parameter
trade-off plots and placement comparison tables; perform failure analysis;
draft final report and clean the codebase. &
Both \\

10 &
Finalize report, figures, README, requirements, and single-command
reproduction workflow; verify results from a clean environment and prepare the
final presentation. \textbf{(Final~due)} &
Both \\

\bottomrule
\end{tabular}
\end{center}

\subsection{Contribution Declaration}

\begin{center}
\begin{tabular}{p{3cm}p{9cm}c}
\toprule
\textbf{Member} & \textbf{Primary Responsibilities} & \textbf{Estimated \%} \\
\midrule

Choo Wen Bin &
Project Lead; experiment design; LoRA implementation; architectural-placement
experiments; experiment integration; statistical analysis; report coordination
and final submission. &
50\% \\

Nai Junn Rong &
Research Lead; literature review; dataset and evaluation pipeline; Prefix
Tuning and full fine-tuning baselines; human factual-consistency study;
reproducibility checks and documentation. &
50\% \\

\bottomrule
\end{tabular}
\end{center}

\section{References}

\begin{enumerate}[label={[\arabic*]},leftmargin=*]

\item
Lewis, M., Liu, Y., Goyal, N., Ghazvininejad, M., Mohamed, A.,
Levy, O., Stoyanov, V., and Zettlemoyer, L., 2020.
``BART: Denoising Sequence-to-Sequence Pre-training for Natural Language
Generation, Translation, and Comprehension.''
\textit{Proceedings of the 58th Annual Meeting of the Association for
Computational Linguistics (ACL)}, pp. 7871--7880.
DOI: \nolinkurl{10.18653/v1/2020.acl-main.703}.

\item
Hu, E. J., Shen, Y., Wallis, P., Allen-Zhu, Z., Li, Y., Wang, S.,
Wang, L., and Chen, W., 2022.
``LoRA: Low-Rank Adaptation of Large Language Models.''
\textit{International Conference on Learning Representations (ICLR)}.
arXiv:2106.09685.

\item
Li, X. L. and Liang, P., 2021.
``Prefix-Tuning: Optimizing Continuous Prompts for Generation.''
\textit{Proceedings of the 59th Annual Meeting of the Association for
Computational Linguistics and the 11th International Joint Conference on
Natural Language Processing (ACL-IJCNLP)}, pp. 4582--4597.
DOI: \nolinkurl{10.18653/v1/2021.acl-long.353}.

\item
Zhang, Q., Chen, M., Bukharin, A., He, P., Cheng, Y., Chen, W.,
and Zhao, T., 2023.
``Adaptive Budget Allocation for Parameter-Efficient Fine-Tuning.''
\textit{International Conference on Learning Representations (ICLR)}.
arXiv:2303.10512.

\item
Veprikov, A., Solodkin, V., Alexander, Z., Savchenko, A.,
and Beznosikov, A., 2026.
``WeightLoRA: Keep Only Necessary Adapters.''
\textit{Proceedings of the 64th Annual Meeting of the Association for
Computational Linguistics (ACL), Volume 1: Long Papers},
pp. 12419--12437.
DOI: \nolinkurl{10.18653/v1/2026.acl-long.566}.

\item
Kryscinski, W., McCann, B., Xiong, C., and Socher, R., 2020.
``Evaluating the Factual Consistency of Abstractive Text Summarization.''
\textit{Proceedings of the 2020 Conference on Empirical Methods in
Natural Language Processing (EMNLP)}, pp. 9332--9346.
DOI: \nolinkurl{10.18653/v1/2020.emnlp-main.750}.

\end{enumerate}

\end{document}
